\scanpage{286}
\mspage{14}

2.7. \deleted{\unsure{8|7}.} Considérons maintenant le foncteur composé

\diagram{\[
\begin{tikzcd}[ampersand replacement=\&]
\mathrm{Ens} \arrow[r,"{I\mapsto I_{\mathrm{Cat}}}"] \& \mathrm{Cat} \arrow[r,"{\mathrm{hot}_W}"] \& \mathrm{Hot}_W,
\end{tikzcd}
\qquad I\longrightarrow \substack{I_{\mathrm{Hot}_W}\\\text{ou } I_W},
\tag*{(2.7.1)}
\]}

où le premier foncteur est celui qui associe à l'ens $I$ la catégorie discrète $I_{\mathrm{Cat}}$ ayant $I$ comme ens d'objets. (Le plus souvent, on identifie $I_{\mathrm{Cat}}$ avec $I$ dans les notations, tout comme pour un ens ordonné $J$, on identifie $J$ et la ``catégorie ordonnée'' définie par $J$. Ici $I$ peut être considéré comme muni de l'ordre discret.) On a d'ailleurs

\[
 I_{\mathrm{Cat}}\simeq I\times e\quad (\deleted{\unsure{\mathrm{i.e.}}}\overset{\mathrm{def}}{=}\text{ somme de }I\text{ copies de }e),
 \tag*{(2.7.2)}
\]

donc quand \deleted{$\mathrm{hot}_W$} $\operatorname{Loc}(4a)$ est satisfait \deleted{\unsure{alors}}\unsure{donc|alors} $\mathrm{hot}_W$ commute aux sommes \unsure{directes|discrètes} on aura un isom canonique

\[
 I_W\simeq I\times e_W. \tag*{(2.7.3.)}
\]

(où $e_W$ est l'objet final de $\mathrm{Hot}_W$, i.e. $e$ considéré comme objet de $\mathrm{Hot}_W$).

\uline{Prop}. 2.8. Supposons \deleted{que} l'axiome \deleted{$\operatorname{Loc}\,\unsure{3|5}$} de connexité $\operatorname{Loc}(5)$ satisfait, d'où un foncteur

\[
 \pi_0:\mathrm{Hot}_W\longrightarrow\mathrm{Ens}. \tag*{(2.8.1)}
\]

Ce foncteur est \uline{adjoint à gauche} du foncteur $I\mapsto I_W$ (2.7.1), i.e. on a une bijection fonctorielle en $(X,I)$ de $\mathrm{Hot}_W\times\mathrm{Ens}$ :

% archive sheet number: 285
